\documentclass[11pt]{article}
\usepackage{amsmath,amssymb,amsthm}
\newtheorem{proposition}{Proposition}
\title{A positive antisymmetric parent has a common child limit}
\author{}\date{}
\begin{document}\maketitle
\paragraph{Status and scope.}
This is a complete written corollary of the exact certificate in
\texttt{positive-paired-root-h16.tex}. It concerns the positive
antisymmetric numerator subfamily only. It does not exclude unrestricted
signed Haar decoders, adaptive pole shapes, or either requested
asymptotic bound. It is not Lean-checked.

\begin{proposition}
Fix $M$ distinct poles $0\leq p_1<\cdots<p_M<1/2$ and signs
$u_i\in\{-1,1\}$. Put
\[
 P(t)=\prod_i(t-p_i),\qquad
 Q_\delta(t)=P(t)P(t-\delta),\qquad
 N_\delta(t)=(t-1/2)R_\delta(t(1-t)),
\]
where $\deg R_\delta\leq M-1$. On the paired poles $p_i,p_i+\delta$
use root masks $u_i,-u_i$ and residues
\[
 c_{i,0}=\frac{N_\delta(p_i)}{u_iQ_\delta'(p_i)},
 \qquad
 c_{i,1}=\frac{N_\delta(p_i+\delta)}{-u_iQ_\delta'(p_i+\delta)}.
\]
Suppose a sequence $\delta\downarrow0$ has all these residues positive.
After multiplying each residue vector by a positive scalar, a subsequence
has a limit of the form $(w,w)$, where $w\geq0$ and $w\ne0$.
\end{proposition}
\begin{proof}
Normalize the coefficient vector of $R_\delta$ to Euclidean norm one;
this is a positive rescaling. Pass to a subsequence on which it converges
to a nonzero polynomial $R_0$ of degree at most $M-1$, and let
$N_0(t)=(t-1/2)R_0(t(1-t))$. Since
\[
 \delta^{-1}Q_\delta'(p_i)\longrightarrow-P'(p_i)^2,
 \qquad
 \delta^{-1}Q_\delta'(p_i+\delta)\longrightarrow P'(p_i)^2,
\]
both $\delta c_{i,0}$ and $\delta c_{i,1}$ tend to
\[
 w_i=-\frac{N_0(p_i)}{u_iP'(p_i)^2}.
\]
Positivity gives $w_i\geq0$. The numbers $p_i(1-p_i)$ are distinct,
and $p_i-1/2\ne0$. If every $w_i$ were zero, $R_0$ would have $M$
distinct zeros, contradicting its degree and normalization.
\end{proof}

\paragraph{Application to the saved sixteen-pole obstruction.}
Fix any sufficiently small positive $\epsilon$ for which the supported
dual in \texttt{positive-paired-root-h16.tex} is positive. Its $17$
supported comparison and nonnegativity rows form a $17$ by $16$
matrix $A$ of rank $16$: the certificate's augmented normalized-dual
matrix $[A^\top;\mathbf1^\top]$ is invertible. If $Aw\geq0$, the
strictly positive left null vector forces $Aw=0$, hence $w=0$.
Thus the two saved root words
\[
 (+,+,+,+,-,-,-,-),\qquad (+,-,+,+,-,-,+,+)
\]
have no common nonzero weak nonnegative residue pool, even when each
has its own bias. Biases have already been eliminated into comparisons
between oppositely labelled queries.

Take these as the two child root words of a thirty-two-pole parent,
with base poles
\[
 p_y=\sum_{j=0}^3 y_j\epsilon^{2^j},\qquad 0\leq y<16,
\]
and newest pairs $p_y,p_y+\delta$. Any fixed parent root word gives
the masks $u_y,-u_y$ in the proposition. Top routing may exchange the
two members of an individual pair, but their limiting residues are
equal, so both children still receive the same vector $w$.

If positive antisymmetric parents realized the required child words
for a sequence $\delta\downarrow0$, scale as in the proposition and
pass to the limit in the bias-eliminated child inequalities. Their
denominators stay bounded away from zero. The common nonzero vector
$w\geq0$ would satisfy all supported weak inequalities, contradicting
$Aw\geq0\Rightarrow w=0$. Therefore, for each such fixed $\epsilon$,
all sufficiently small newest gaps $\delta$ are infeasible in this
subfamily. Neither diverging coefficients nor a boundary residue pool
avoids this conclusion.

\paragraph{Quantifier caution.}
The conclusion is $\forall\epsilon$ sufficiently small,
$\exists\delta_0(\epsilon)>0$, followed by infeasibility for
$0<\delta<\delta_0(\epsilon)$. No estimate on $\delta_0$ is supplied.
In particular this argument alone does not assert infeasibility along
the specific curve $\delta=\epsilon^{16}$.
\end{document}
